\subsection{Resolvent topologies and continuity of the functional calculus}

In this issue, E denotes a complex Hilbert space. Recall that we denote by \(\mathcal{A}(\mathrm{E})\) the set of partial self-adjoint operators on E. We will extend to \(\mathcal{A}(\mathrm{E})\) the continuity properties of item 8 of IV, p. 194.

DEFINITION 6. --- Let T be a locally convex topology on \(\mathcal{L}(\mathrm{E})\) . The T-resolvent topology on \(\mathcal{A}(\mathrm{E})\) is the coarsest topology such that the maps \(u \mapsto \mathrm{R}(u, \lambda)\) from \(\mathcal{A}(\mathrm{E})\) into \(\mathcal{L}(\mathrm{E})\) endowed with the topology T are continuous for every \(\lambda \in \mathbf{C} - \mathbf{R}\) .

PROPOSITION 11. --- Let T be a locally convex topology on \(\mathcal{L}(\mathrm{E})\) which is coarser than the Banach space topology of \(\mathcal{L}(\mathrm{E})\) .

\begin{enumerate}
\def\labelenumi{\alph{enumi})}
\item
  Let \(f \in \mathcal{C}_0(\mathbf{R})\) . For every sequence \((u_n)_{n \in \mathbf{N}}\) in \(\mathcal{A}(\mathrm{E})\) which converges to u for the topology \(\mathcal{T}\) -resolvent, the sequence \((f(u_n))_{n \in \mathbf{N}}\) converges to \(f(u)\) in the space \(\mathcal{L}(\mathrm{E})\) endowed with the topology \(\mathcal{T}\) ;
\item
  Suppose that everything \(u \in \mathcal{A}(\mathrm{E})\) admits a countable fundamental system of neighborhoods for the topology \(\mathcal{T}\) -resolvent. The map of \(\mathcal{C}_0(\mathbf{R}) \times \mathcal{A}(\mathrm{E})\) in the space \(\mathcal{L}(\mathrm{E})\) endowed with the topology \(\mathcal{T}\) defined by \((f, u) \mapsto f(u)\) is continuous.
\end{enumerate}

Let us prove \(a\) ). The topology \(\mathcal{T}\) on \(\mathcal{L}(\mathrm{E})\) is the upper bound of the topologies defined by the semi-norms continuous for the topology \(\mathcal{T}\) (EVT, II, p. 26, cor. and following remark). It therefore suffices to prove that for every seminorm \(p\) on \(\mathcal{L}(\mathrm{E})\) which is continuous for the topology \(\mathcal{T}\) the sequence \(p(f(u_n) - f(u))\) converges to 0 (TG, I, p. 51, prop. 10).

Let p be such a seminorm. Denote by \(\mathcal{L}(\mathrm{E})_{p}\) the separated Banach space completion of the space \(\mathcal{L}(\mathrm{E})\) equipped with the seminorm p, and denote by \(\varpi\) the canonical map from \(\mathcal{L}(\mathrm{E})\) into \(\mathcal{L}(\mathrm{E})_{p}\) ; it is continuous, and one has \(p(u) = \|\varpi(u)\|_{p}\) for all \(u \in \mathcal{L}(\mathrm{E})\) , where the norm is that of the space \(\mathcal{L}(\mathrm{E})_{p}\) .

Since T is coarser than the Banach space topology of \(\mathcal{L}(\mathrm{E})\) , there exists \(c \geqslant 0\) such that \(p(u) \leqslant c \|u\|\) for all \(u \in \mathcal{L}(\mathrm{E})\) .

Suppose first that \(f \in \mathcal{D}(\mathbf{R})\) Let \(\tilde{f}\) an almost analytic extension of \(f\) . Denote by \(\mu\) the Lebesgue measure on \(\mathbf{C}\) . By the Helffer--Sjöstrand formula (th. 2 of IV, p. 284), one has

\[
\varpi (f (u _ {n})) = - \frac {1}{\pi} \int_ {\mathbf {C}} \frac {\partial \widetilde {f}}{\partial \bar {z}} (\lambda) \varpi (\mathrm{R} (u _ {n}, \lambda)) d \mu (\lambda)\tag{14}
\]

for all \(n \in \mathbf{N}\) .

Let \(\lambda \in \mathbf{C} - \mathbf{R}\) By hypothesis, the sequence of resolvents \(\mathrm{R}(u_n,\lambda)\) converges to \(\mathrm{R}(u,\lambda)\) in \(\mathcal{L}(\mathrm{E})\) endowed with the topology \(\mathcal{T}\) , hence the sequence \((\varpi (\mathrm{R}(u_n,\lambda)))_n\) converges to \(\varpi (\mathrm{R}(u,\lambda))\) to the Banach space \(\mathcal{L}(\mathrm{E})_p\) .

For every \(n \in \mathbf{N}\) , we have

\[
\left\| \frac {\partial \widetilde {f}}{\partial \bar {z}} (\lambda) \mathrm{R} (u _ {n}, \lambda) \right\| \leqslant \left| \frac {1}{\mathcal{I} (\lambda)} \frac {\partial \widetilde {f}}{\partial \bar {z}} (\lambda) \right|
\]

(prop. 17 of IV, p. 248) therefore

\[
\left\| \frac {\partial \widetilde {f}}{\partial \overline {{z}}} (\lambda) \varpi (\mathrm{R} (u _ {n}, \lambda)) \right\| _ {p} \leqslant c \left| \frac {1}{\mathcal {I} (\lambda)} \frac {\partial \widetilde {f}}{\partial \overline {{z}}} (\lambda) \right|.
\]

The estimate (13) of IV, p. 281 shows that the right-hand side of this inequality is a bounded function for \(\lambda \in \mathbf{C} - \mathbf{R}\) ; it is integrable on \(\mathbf{C}\) since \(\widetilde{f}\) has compact support. One deduces from Lebesgue's theorem (INT, IV, p. 137, § 3, n° 7, th. 6) and from the Helffer--Sjöstrand formula applied to \(f\) we deduce that \(\varpi(f(u_n))\) converges to \(\varpi(f(u))\) , hence \(p(f(u_n) - f(u))\) tends to 0.

Suppose now that \(f \in \mathcal{C}_0(\mathbf{R})\) Let \(\varepsilon > 0\) . There exists a function \(f_{\varepsilon}\) in \(\mathcal{D}(\mathbf{R})\) such that \(\| f_{\varepsilon} - f \|_{\infty} \leqslant \varepsilon\) (cf. prop. 4, a) of IV, p. 202 and INT, III, p. 45, § 1, n° 2, prop. 3). According to prop. 5, c) of IV, p. 275, one then has \(\| f(u) - f_{\varepsilon}(u) \| \leqslant \varepsilon\) and \(\| f(u_n) - f_{\varepsilon}(u_n) \| \leqslant \varepsilon\) for all \(n \in \mathbf{N}\) . Consequently, it follows that

\[
p (f (u _ {n}) - f (u)) \leqslant 2 c \varepsilon + p (f _ {\varepsilon} (u _ {n}) - f _ {\varepsilon} (u))
\]

for all \(n \in \mathbf{N}\) . Since \(f_{\varepsilon} \in \mathcal{D}(\mathbf{R})\) , the sequence \((p(f_{\varepsilon}(u_n) - f_{\varepsilon}(u)))_{n \in \mathbf{N}}\) converges to 0 by the preceding case, and therefore \(p(f(u_n) - f(u))\) converges to 0. This concludes the proof of a).

Let us prove assertion \(b\) ). Under the hypothesis concerning \(\mathcal{T}\) , it follows from TG, IX, p. 17, prop. 10 and the following remark, and from assertion \(a\) ), that the map \(u \mapsto f(u)\) of \(\mathcal{A}(\mathrm{E})\) in \(\mathcal{L}(\mathrm{E})\) is continuous when \(f \in \mathcal{C}_0(\mathbf{R})\) . Since the maps \(f \mapsto f(u)\) of \(\mathcal{C}_0(\mathbf{R})\) in \(\mathcal{L}(\mathrm{E})\) are continuous with norm \(\leqslant 1\) for all \(u \in \mathcal{A}(\mathrm{E})\) (prop. 5, \(c\) ) of IV, p. 275), assertion \(b\) ) follows from TG, X, p. 13, cor. 3.

Examples. --- 1) Let \(\mathscr{S}\) be a set of bounded subsets of E. The \(\mathscr{S}\)-topology on \(\mathcal{L}(\mathrm{E})\) (EVT, III, p. 13) is a locally convex topology less fine than the Banach space topology of \(\mathcal{L}(\mathrm{E})\) .

\begin{enumerate}
\def\labelenumi{\arabic{enumi})}
\setcounter{enumi}{1}
\tightlist
\item
  Let \(T_{b}\) the Banach space topology of \(\mathcal{L}(\mathrm{E})\) . For the topology \(T_{b}\) -resolvent, every \(u \in \mathcal{A}(\mathrm{E})\) admits a countable fundamental system of neighborhoods.
\end{enumerate}

Indeed, it suffices to show that for every \(\lambda \in \mathbf{C} - \mathbf{R}\) and every \(\varepsilon > 0\) , there exist \(\lambda' \in \mathbf{Q} + i\mathbf{Q}^*\) and an integer \(n \geqslant 1\) such that every partial operator \(v \in \mathcal{A}(\mathrm{E})\) satisfying \(\| \mathrm{R}(v, \lambda') - \mathrm{R}(u, \lambda')\| < 1/n\) also satisfies condition \(\|R(v,\lambda)-R(u,\lambda)\|<\varepsilon\) ; by writing

\[
\begin{array}{c} \| \mathrm{R} (v, \lambda) - \mathrm{R} (u, \lambda) \| \leqslant \| \mathrm{R} (v, \lambda) - \mathrm{R} (v, \lambda^ {\prime}) \| \\ + \| \mathrm{R} (v, \lambda^ {\prime}) - \mathrm{R} (u, \lambda^ {\prime}) \| + \| \mathrm{R} (u, \lambda^ {\prime}) - \mathrm{R} (u, \lambda) \|, \end{array}
\]

this property follows from formula (8) of IV, p. 245, from the estimates of proposition 17 of IV, p. 248, and from the fact that \(\mathbf{Q} + i\mathbf{Q}^*\) is everywhere dense in \(\mathbf{C} - \mathbf{R}\) .

The conclusion of the proposition is not valid in general if \(\mathcal{C}_{0}(\mathbf{R})\) is replaced by \(\mathcal{C}_{b}(\mathbf{R})\) (exercise 29 of IV, p. 366, e)). Nevertheless, one has the following result.

COROLLARY. --- Let \(\mathcal{T}_{s}\) the topology of simple convergence on \(\mathcal{L}(\mathrm{E})\) Let \(f \in \mathcal{C}_{b}(\mathbf{R})\) . For every sequence \((u_{n})_{n \in \mathbf{N}}\) in \(\mathcal{A}(\mathrm{E})\) which converges to \(u\) for the topology \(\mathcal{T}_{s}\) -resolvent, the sequence \((f(u_{n}))_{n \in \mathbf{N}}\) converges to \(f(u)\) in \(\mathcal{L}(\mathrm{E})\) endowed with the topology \(\mathcal{T}_{s}\) .

Let \((u_{n})\) be a sequence in \(\mathcal{A}(\mathrm{E})\) which converges to u for the topology \(T_{s}\) -resolvent. Let \(x \in E\) and let \(\varepsilon > 0\) .

For every integer \(N \in \mathbf{N}\) , let \(\varphi_{N} \in \mathcal{K}(\mathbf{R})\) be a function with support contained in \([-(N+1), N+1]\) such that \(0 \leqslant \varphi_{N} \leqslant 1\) and \(\varphi_{N}(t) = 1\) for all \(t \in [-N, N]\) . The functions \(\varphi_{N}\) converge pointwise to 1 as N tends to infinity, and satisfy \(|\varphi_{N}| \leqslant 1\) ; hence prop. 6 of IV, p. 276 implies that there exists \(N \in \mathbf{N}\) such that \(\|\varphi_{N}(u)x - x\| \leqslant \varepsilon\) . Set \(f_{N} = f\varphi_{N}\) .

For every \(n\in \mathbf{N}\) , we have

\[
\begin{array}{r l} \| f (u _ {n}) x - f (u) x \| \leqslant & \| f (u _ {n}) x - f _ {\mathrm{N}} (u _ {n}) x \| + \\ & \| f _ {\mathrm{N}} (u _ {n}) x - f _ {\mathrm{N}} (u) x \| + \| f _ {\mathrm{N}} (u) x - f (u) x \|. \end{array} \tag {15}
\]

For every \(v\in \mathcal{A}(\mathrm{E})\) , we have

\[
\begin{array}{r l} & {\| f (v) x - f _ {\mathrm{N}} (v) x \| = \| (f (1 - \varphi_ {\mathrm{N}})) (v) x \|} \\ & {\qquad \leqslant \| f (v) \|   \| x - \varphi_ {\mathrm{N}} (v) x \| \leqslant \| f \| _ {\infty}   \| x - \varphi_ {\mathrm{N}} (v) x \|} \end{array}
\]

(prop. 5, c) of IV, p. 275). Since \(\varphi_{\mathrm{N}} \in \mathcal{C}_0(\mathbf{R})\) , assertion \(a\) ) of prop. 11 applied to the topology \(\mathcal{T}_s\) -resolvent implies that \(\varphi_{\mathrm{N}}(u_n)x\) converges to \(\varphi_{\mathrm{N}}(u)x\) as \(n \to +\infty\) . Consequently, for every \(n\) sufficiently large, one has

\[
\left\| x - \varphi_ {\mathrm{N}} (u _ {n}) x \right\| \leqslant \left\| x - \varphi_ {\mathrm{N}} (u) x \right\| + \varepsilon \| x \| \leqslant (1 + \| x \|) \varepsilon .
\]

Inequality (15) then becomes

\[
\| f (u _ {n}) x - f (u) x \| \leqslant \| f \| _ {\infty} (2 + \| x \|) \varepsilon + \| f _ {\mathrm{N}} (u _ {n}) x - f _ {\mathrm{N}} (u) x \|
\]

for all sufficiently large n. One concludes the proof by means of assertion a) of loc. cit., applied to the function \(f_{\mathrm{N}} \in \mathcal{C}_{0}(\mathbf{R})\) and to the topology \(T_{s}\) -resolvent.

